\section{Проектирование решателя задач}

\subsection{Определение пользователей решателя задач}
% Определяются категории пользователей и подсистем, взаимодействующих с решателем задач. Для каждой категории указываются роль, цели, уровень подготовки, полномочия и доступные функции.

% Например: эксперт предметной области, оператор, конечный пользователь, администратор базы знаний, внешний программный сервис.

% Для каждой категории пользователей формулируются цели взаимодействия с решателем задач. Указываются ожидаемые результаты: получение решения, диагностика, классификация, рекомендация, конфигурация, план действий, объяснение результата или проверка условий.

% Следует отделить пользовательские цели от внутренних функций системы.


Взаимодействие с решателем задач осуществляется через несколько категорий участников, различающихся ролью, целями, уровнем подготовки и доступными функциями. Ниже описаны пользовательские категории и внешние подсистемы, а также проведено разграничение между пользовательскими целями и внутренними функциями системы.

\subsubsection{Конечный пользователь}

\textbf{Роль.} Исследователь, аналитик или инженер, накапливающий личную базу знаний из неструктурированных документов.

\textbf{Цели взаимодействия.}
\begin{itemize}
    \item пополнение базы знаний новыми документами;
    \item получение ответа на вопрос с обоснованием (путём рассуждения и ссылками на первоисточники);
    \item навигация по иерархии «сообщество $\rightarrow$ домен $\rightarrow$ концепция $\rightarrow$ карточка»;
    \item визуальный обзор структуры графа знаний.
\end{itemize}

\textbf{Уровень подготовки.} Базовый --- не требует знания внутренних механизмов системы (эмбеддинги, алгоритмы обхода графа, работа LLM).

\textbf{Полномочия.} Чтение и запись в пределах собственной базы знаний; просмотр карточек и оригиналов; удаление загруженных документов с каскадным обновлением графа.

\textbf{Доступные функции.} Загрузка документов (UC-1), чат с вопросами (UC-2), навигация по графу (UC-2), визуализация (UC-4).

\subsubsection{Администратор базы знаний}

\textbf{Роль.} Продвинутый пользователь, отвечающий за качество и целостность накопленных знаний.

\textbf{Цели взаимодействия.}
\begin{itemize}
    \item контроль корректности канонического реестра концепций;
    \item разрешение конфликтов нормализации (слияние дубликатов, разделение омонимов);
    \item инициирование пересчёта кластеризации и перестройки иерархии;
    \item мониторинг ошибок конвейера обработки.
\end{itemize}

\textbf{Уровень подготовки.} Продвинутый --- понимание модели графа, принципов нормализации и оркестрации агентов.

\textbf{Полномочия.} Редактирование реестра концепций, слияние и удаление узлов, просмотр журналов обработки, ручной запуск агента-кластеризатора.

\textbf{Доступные функции.} Модерация реестра, просмотр логов, пересчёт кластеризации.

\subsubsection{Внешние подсистемы}

Помимо пользователей, с решателем взаимодействуют внешние программные сервисы, не являющиеся пользователями в прямом смысле, но участвующие в обработке данных.

\begin{itemize}
    \item \textbf{LLM-провайдер} --- внешняя или локальная языковая модель, вызываемая агентами-аналитиком, нормализатором и ответчиком. Вход: промпт и контекст; выход: структурированный JSON или текст.
    \item \textbf{Векторный индекс} (FAISS / HNSW / Qdrant) --- хранилище эмбеддингов узлов и запросов. Вход: вектор запроса; выход: top-$K$ ближайших узлов.
    \item \textbf{Объектное хранилище} (S3-совместимое: Neon Object Storage, Cloudflare R2) --- хранение оригиналов документов. Вход: файл; выход: ссылка на скачивание.
    \item \textbf{СУБД} --- хранение концепций, карточек, связей, реестра, доменов и сообществ.
\end{itemize}

\textbf{Разграничение пользовательских целей и внутренних функций}

Важно отделять цели, ради которых пользователь обращается к системе, от внутренних функций, которые система выполняет для их достижения (таблица~\ref{table:goals_functions}).

\begin{table}[H]
\centering
\caption{Пользовательские цели и внутренние функции системы}
\label{table:goals_functions}
\begin{tabularx}{\textwidth}{|>{\RaggedRight\arraybackslash}X|>{\RaggedRight\arraybackslash}X|}
\hline
\textbf{Пользовательская цель} & \textbf{Внутренние функции системы} \\ \hline
Получить ответ на вопрос & эмбеддинг запроса, векторный поиск точек входа, локальный обход графа, ранжирование подграфа, генерация ответа LLM, построение пути рассуждения \\ \hline
Добавить документ в базу & парсинг файла, извлечение концепций и summary, нормализация концепций, расчёт весов рёбер, обновление карточек \\ \hline
Увидеть структуру знаний & кластеризация графа, генерация описаний сообществ и доменов, отдача данных для визуализации \\ \hline
Найти первоисточник & хранение оригинала в объектном хранилище, привязка ссылки к карточке \\ \hline
\end{tabularx}
\end{table}

Такое разграничение важно для проектирования: пользователь формулирует цель на естественном языке, а система реализует её через цепочку внутренних агентов и алгоритмов, скрытых от пользователя.




\subsection{Сценарии использования решателя задач}
% Описываются сценарии использования (use cases), отражающие взаимодействие пользователя или внешней системы с решателем задач.
% Необходимо выделить типичные, граничные и ошибочные сценарии.


Для детального проектирования взаимодействия пользователя с системой разработаны сценарии использования (use cases). Они разделены на типичные (основные сценарии работы), граничные (работа в условиях предельных нагрузок или неполных данных) и ошибочные (обработка сбоев).

\textbf{Типичные сценарии.}

\vspace{0.5em}
\textbf{1 Пополнение базы знаний (Загрузка нового документа).}

Актором выступает конечный пользователь, работающий через веб-интерфейс. Предусловием является авторизация пользователя в системе. В основном потоке пользователь загружает файл (PDF, Markdown, TXT) через экран загрузки, при этом фронтенд отображает индикатор прогресса обработки. Затем агент-аналитик извлекает текст, определяет домен, генерирует summary, ключевые тезисы и список концепций. Далее агент-нормализатор сверяет концепции с реестром, разрешает синонимы и создает новые узлы при необходимости. После этого агент-связыватель вычисляет веса рёбер и связывает новую карточку с ранее загруженными документами по общим концепциям. В завершение система возвращает результат: карточку знания, количество созданных концепций и добавленных связей. Постусловием является обновление графа знаний, сохранение оригинала файла в объектном хранилище и доступность карточки для навигации.

\vspace{0.5em}
\textbf{2 Поиск ответа на вопрос (Диалоговый режим).}

Актором выступает конечный пользователь. Предусловием является наличие в базе знаний хотя бы одного обработанного документа. В основном потоке пользователь вводит вопрос на естественном языке в чат-интерфейсе. Агент-ответчик выполняет эмбеддинг запроса и векторный поиск top-$K$ концепций (точек входа). Выполняется локальный обход графа (1--2 хопа) с использованием beam search. Найденный подграф ранжируется по весам рёбер и близости к запросу. Языковая модель генерирует ответ на основе отфильтрованного подграфа. В конце фронтенд отображает ответ, путь рассуждения (вопрос $\rightarrow$ концепция $\rightarrow$ первоисточник) и ссылки на оригиналы.

\vspace{0.5em}
\textbf{Граничные сценарии.}

\vspace{0.5em}
\textbf{3 Обработка документа большого объёма.}

Если документ превышает контекстное окно LLM, агент-аналитик разбивает его на логические фрагменты (чанки), извлекает концепции из каждого фрагмента, а затем агрегирует результат в единую карточку знания, избегая дублирования концепций.

\vspace{0.5em}
\textbf{4 Неоднозначный или абстрактный запрос.}

Если векторный поиск не даёт явных точек входа с высокой степенью уверенности, агент-ответчик применяет механизм расширения запроса (query expansion) или предлагает пользователю уточняющий вопрос, отображая в интерфейсе смежные концепции из верхних уровней иерархии.

\vspace{0.5em}
\textbf{Ошибочные сценарии.}

\vspace{0.5em}
\textbf{5 Некорректный формат или повреждение файла.}

Если при загрузке файла парсер выдаёт ошибку, система изолирует сбой (не прерывая работу всего конвейера), уведомляет пользователя через интерфейс о невозможности прочтения файла и предлагает загрузить документ в другом формате.

\vspace{0.5em}
\textbf{6 Сбой внешней языковой модели (Timeout / API Error).}

Если LLM-провайдер недоступен в момент загрузки документа, система сохраняет оригинал файла и ставит задачу в очередь на повторную обработку. Если сбой происходит при ответе на вопрос, агент-ответчик возвращает пользователю только релевантный подграф (найденные концепции и карточки) без сгенерированного текста, сохраняя ценность структуры базы знаний.

\vspace{0.5em}
\textbf{7 Запрос к пустой базе знаний.}

Если пользователь пытается задать вопрос до загрузки документов, фронтенд перехватывает состояние пустого графа и отображает приглашение загрузить первый документ вместо чата.


\subsection{Архитектура решателя задач}
% Представляется общая архитектура решателя задач и его место в составе интеллектуальной системы. Показываются связи с базой знаний, будущим пользовательским интерфейсом, внешними сервисами и источниками данных.
%
% Рекомендуется привести компонентную или структурную диаграмму. На диаграмме следует выделить основные компоненты и направления обмена данными между ними.
%
% Например:
% - модуль приёма и предварительной обработки запроса;
% - модуль проверки и нормализации входных данных;
% - модуль доступа к базе знаний;
% - машина вывода, поиска или иной механизм решения задач;
% - модуль проверки ограничений;
% - модуль формирования результата;
% - модуль формирования объяснения;
% - модуль журналирования и диагностики.
%
% Для каждого компонента указываются ответственность, входные данные,
% выходные данные и взаимодействие с другими компонентами.

Разрабатываемый решатель задач представляет собой распределённую интеллектуальную систему. В индивидуальной части работы детально проектируется подсистема взаимодействия с пользователем: клиентское веб-приложение, визуализация графа, состояния загрузки документов, просмотр карточек и представление ответа с источниками. Агентное ядро, база знаний и серверная инфраструктура рассматриваются как внешние компоненты, доступные через согласованный REST API.

Архитектура разрабатываемой подсистемы включает следующие ключевые компоненты, обеспечивающие полный цикл работы пользователя с базой знаний.

\vspace{0.5em}
\textbf{1 Модуль пользовательского интерфейса (Frontend).}

Данный компонент реализуется на базе фреймворка React и отвечает за всё визуальное взаимодействие с пользователем. Входными данными являются действия пользователя (выбор файла, ввод текста вопроса). Выходными данными — отображение прогресса обработки документа, навигация по иерархии «сообщество $\rightarrow$ домен $\rightarrow$ концепция», просмотр карточек знаний и чат-интерфейс. Модуль также отвечает за интерактивную визуализацию графа знаний и отображение ссылок на оригинальные документы.

\vspace{0.5em}
\textbf{2 Модуль приёма и предварительной обработки запроса.}

Этот компонент связывает фронтенд с ядром системы. Клиентская часть проверяет расширение файла, количество файлов в одной отправке и непустой текст вопроса, после чего передаёт данные серверу. Серверная часть выполняет окончательную проверку и инициирует соответствующий сценарий агентного ядра.

\vspace{0.5em}
\textbf{3 Агент-аналитик (внешний компонент).}

Интеллектуальный компонент отвечает за первичную обработку загруженного документа и возвращает структурированную карточку знания. Для клиентской подсистемы существенен только контракт результата: аннотация, ключевые тезисы, сообщество, домен, концепции и ссылка на оригинал. Внутренняя реализация агента не относится к индивидуальной части работы.

\vspace{0.5em}
\textbf{4 Машина поиска (внешний агент-ответчик).}

Компонент принимает вопрос пользователя и возвращает текстовый ответ, список подтверждающих карточек и упорядоченные пути по концепциям. Клиентская подсистема не выполняет векторный поиск или обход графа, а проверяет состояние формы, отправляет запрос и визуально разделяет ответ, источники и доказательные пути.

\vspace{0.5em}
\textbf{5 Модуль проверки ограничений и ранжирования.}

Работает в связке с машиной поиска и отвечает за отсечение нерелевантных данных. Входные данные — локальный подграф. Компонент применяет ограничения по глубине обхода, ширине луча и весам рёбер, а также ранжирует узлы по близости к вопросу и графовой центральности. Выходные данные — топ-$N$ узлов и связанных карточек, которые передаются в языковую модель.

\vspace{0.5em}
\textbf{6 Модуль формирования результата и объяснения.}

Отвечает за генерацию финального ответа и обеспечение свойства traceability (прослеживаемости). Входные данные — вопрос пользователя и отфильтрованный топ-$N$ подграф. Выходные данные — текстовый ответ, сгенерированный LLM, и структурированный путь рассуждения (например, «вопрос $\rightarrow$ концепция A $\rightarrow$ карточка B $\rightarrow$ оригинал документа C»). Фронтенд использует эти данные для отрисовки цепочки, по которой пользователь может кликнуть и перейти к первоисточнику.

\vspace{0.5em}
\textbf{7 Модуль журналирования и диагностики.}

Служебный компонент, фиксирующий события на стороне разрабатываемой подсистемы. Входные данные — события от агентов и интерфейса (ошибки загрузки, таймауты LLM). Выходные данные — логи, которые используются для отладки конвейера и анализа узких мест.

\vspace{0.5em}
Взаимодействие клиентского приложения с базой знаний и агентами происходит только через REST API. Такая граница позволяет независимо изменять серверные алгоритмы при сохранении типизированных форматов документа, карточки, графа и ответа.

Компонентная схема индивидуально проектируемой подсистемы приведена на рисунке~\ref{fig:frontend_architecture_design}.

\begin{figure}[h!]
\centering
\begin{tabular}{c}
\fbox{\parbox{0.78\textwidth}{\centering Пользователь: загрузка, навигация по графу, вопрос}}\\
$\downarrow$\\
\fbox{\parbox{0.78\textwidth}{\centering React-приложение: Library, Graph, Ask, Card modal}}\\
$\downarrow$\\
\fbox{\parbox{0.78\textwidth}{\centering API-клиент: авторизация, FormData, JSON, загрузка оригинала}}\\
$\downarrow$\\
\fbox{\parbox{0.78\textwidth}{\centering REST API: документы, задания, карточки, граф, ответы}}\\
$\downarrow$\\
\fbox{\parbox{0.78\textwidth}{\centering Агентное ядро и хранилища как внешние компоненты}}
\end{tabular}
\caption{Компонентная архитектура клиентской подсистемы}
\label{fig:frontend_architecture_design}
\end{figure}

\subsection{Описание процесса решения задач}
% Описывается общий процесс решения задачи: получение запроса, проверка его корректности, интерпретация входных данных, обращение к базе знаний, выполнение вывода или поиска, проверка ограничений, формирование решения, объяснения и ответа.
%
% Следует показать альтернативные ветви процесса: недостаточность данных, противоречивость знаний, отсутствие решения, превышение ограничений поиска и необходимость уточнения запроса.

Процесс решения задачи поиска и генерации ответа представляет собой последовательное взаимодействие клиентского приложения с внешним агентом-ответчиком. Фронтенд управляет вводом, состояниями ожидания и представлением результата, но не выполняет интеллектуальный поиск самостоятельно.

\vspace{0.5em}
\textbf{1 Получение и валидация запроса.}

Процесс начинается с того, что пользователь вводит вопрос на естественном языке в чат-интерфейсе (фронтенд). Модуль валидации проверяет запрос на корректность: он не должен быть пустым или превышать допустимую длину. Если запрос корректен, он передается в API, откуда попадает к агенту-ответчику.

\vspace{0.5em}
\textbf{2 Интерпретация входных данных и обращение к базе знаний.}

Агент-ответчик выполняет эмбеддинг (векторизацию) текста запроса. Затем он обращается к базе знаний через векторный индекс (HNSW). Происходит поиск top-$K$ концепций, наиболее близких к запросу по косинусному расстоянию. Эти концепции становятся «точками входа» в граф.

\vspace{0.5em}
\textbf{3 Выполнение поиска и проверка ограничений.}

От найденных точек входа выполняется локальный обход графа на глубину 1–2 хопа. На этом этапе вступают в силу ограничения поиска: используется алгоритм beam search (ширина луча), который отсекает малозначимые ветви. Вычислительная сложность ограничивается $O(\log N)$ за счет отказа от полного обхода. Найденные узлы и связи ранжируются по весу рёбер и близости к исходному запросу.

\vspace{0.5em}
\textbf{4 Формирование решения и объяснения.}

Отфильтрованный подграф (топ-$N$ узлов) вместе с текстами привязанных к ним карточек передается в большую языковую модель (LLM). Модель генерирует связный текстовый ответ. Одновременно с этим формируется объяснительный путь (traceability), который показывает, через какие концепции и документы прошла система, чтобы дать этот ответ. Фронтенд отображает ответ и цепочку рассуждения со ссылками на первоисточники.

\vspace{0.5em}
\textbf{5 Альтернативные ветви процесса.}

В зависимости от состояния базы знаний и характера запроса, процесс может отклоняться от основного сценария. Если база знаний пуста или векторный поиск не находит ни одной концепции с достаточной степенью уверенности, система не вызывает LLM, а возвращает пользователю сообщение о недостаточности данных с предложением загрузить новые документы. Если запрос сформулирован слишком абстрактно или неоднозначно, агент-ответчик применяет механизм расширения запроса (query expansion) или формирует уточняющий вопрос, предлагая пользователю выбрать из смежных концепций верхнего уровня. В случае, если в графе обнаруживаются противоречивые сведения (например, разные карточки трактуют концепцию по-разному), LLM явно указывает на это в тексте ответа, а путь рассуждения отображает обе ветви. Если количество найденных связей превышает лимиты контекстного окна, срабатывают ограничения поиска: система отсекает менее релевантные узлы и явно уведомляет об этом пользователя в интерфейсе.

\subsection{Алгоритмы решения задач}
% Детально описываются алгоритмы, реализующие выбранные методы решения задач. Указываются этапы работы, входные и выходные данные, условия переходов, правила выбора действий, критерии остановки и порядок обращения к базе знаний.
%
% Алгоритм рекомендуется представить в виде псевдокода, блок-схемы, диаграммы деятельности или диаграммы последовательности.

% Анализируется вычислительная сложность каждого проектируемого алгоритма. Рассматривается влияние объёма базы знаний, числа правил, глубины поиска, количества ограничений и сложности запроса на время выполнения и потребление вычислительных ресурсов.
%
% Формулируются проектные ограничения и меры по их снижению: ограничение глубины поиска, эвристики, кэширование, индексация знаний, предварительная фильтрация фактов или декомпозиция задачи.
В основе проектируемой клиентской подсистемы лежат алгоритмы координации состояния интерфейса, раскладки графа и отображения асинхронных результатов. Алгоритмы извлечения знаний и двухступенчатого поиска выполняются внешним агентным ядром и используются через API-контракты.

\vspace{0.5em}
\textbf{1 Алгоритм координации состояния документа.}

Входными данными являются выбранные пользователем файлы. Клиент отбрасывает неподдерживаемые расширения, ограничивает пакет десятью файлами и отправляет допустимые данные. Созданные задания сразу добавляются в список документов, после чего клиент периодически получает состояния только активных заданий. Опрос прекращается, когда все задания переходят в конечное состояние или пользователь покидает экран. Для завершённого задания становятся доступны карточка и оригинал, для ошибочного --- повторная обработка, а для выполняемого блокируется удаление.

\vspace{0.5em}
\textbf{2 Алгоритм иерархической раскладки графа.}

Входными данными являются массивы сообществ, доменов, концепций и рёбер. Для каждого сообщества выделяется прямоугольная область. Домены размещаются в сетке внутри области сообщества, а концепции --- кольцами вокруг соответствующего домена. Координаты всех доменов и концепций сохраняются в карте, после чего по ней строятся доменные и межконцептуальные рёбра. Выходом является детерминированная SVG-композиция, пригодная для масштабирования, перемещения и интерактивного выделения соседей.

\vspace{0.5em}
\textbf{3 Алгоритм представления ответа.}

После проверки непустого вопроса клиент блокирует повторную отправку и отображает состояние ожидания. Успешный ответ разделяется на три визуальных блока: текст, список источников и пути по концепциям. Для каждого источника сохраняется действие скачивания оригинала; для пути показываются последовательность концепций и числовая релевантность. При ошибке предыдущий экран остаётся доступным, а понятное сообщение выводится рядом с формой.

\vspace{0.5em}
\textbf{4 Анализ вычислительной сложности и проектные ограничения.}

Построение координат выполняется за $O(C+D+V)$, где $C$ --- число сообществ, $D$ --- доменов, $V$ --- концепций. Отрисовка межконцептуальных связей требует $O(E)$ операций, где $E$ --- число рёбер. Выделение соседей выбранной концепции также требует одного просмотра массива рёбер, то есть $O(E)$. Для прототипа этого достаточно, однако при значительном росте графа необходимы виртуализация, пространственный индекс или серверная выдача ограниченного подграфа. Отображение карточек концепции выполняется по запросу и кэшируется в состоянии компонента, что исключает повторную загрузку при последующих наведениях.

\subsection{Выводы}

Спроектирована клиентская подсистема, которая отделена от агентного ядра типизированными REST-контрактами. Определены состояния авторизации, библиотеки, графа, карточки и диалога, а также правила представления ожидания, ошибок и конечных результатов.

Разработаны алгоритмы координации статусов документов, иерархической SVG-раскладки и представления ответа с источниками. Агентные алгоритмы рассмотрены только как внешние поставщики структурированных данных, что соответствует индивидуальной зоне ответственности автора.
